\begin{papercard}{nowell2017}
\cf{Stages} \stagetag{C} (also \stagetag{I} \stagetag{R}). A practical map of how to make Braun \& Clarke's six-phase thematic analysis \emph{trustworthy} in Lincoln \& Guba's terms (credibility, transferability, dependability, confirmability). Table 1 lists concrete trustworthiness actions for each phase, and each phase is illustrated by a mixed-methods case study. It is a checklist source for audit-trail and documentation practices.

\cf{Aim and gap} Qualitative analysis is often under-reported, and the researcher is the analytic instrument, so rigor needs disclosure with enough detail for readers to judge credibility \pc{nowell2017}{2}{Qualitative researchers can demonstrate how data analysis has been conducted through recording, systematizing, and disclosing the methods}. Thematic analysis (TA) is flexible and accessible, but its flexibility can cause inconsistency and incoherent themes unless an epistemological position is stated \pc{nowell2017}{2}{While thematic analysis is flexible, this flexibility can lead to inconsistency and a lack of coherence}. The authors treat TA as a method in its own right.

\cf{Trustworthiness criteria} \emph{Credibility}: fit between respondents' views and the researcher's representation; techniques include prolonged engagement, triangulation of data and researchers, peer debriefing, referential adequacy and member checking. \emph{Transferability}: thick description for case-to-case transfer. \emph{Dependability}: a logical, traceable, documented process that can be audited. \emph{Confirmability}: interpretations clearly derived from the data, with reasons for choices recorded \pc{nowell2017}{3}{Confirmability is concerned with establishing that the researcher’s interpretations and findings are clearly derived from the data}. \emph{Audit trail}: raw data, field notes, transcripts and a \emph{reflexive journal} that records logistics, methodological decisions and rationales, and the researcher's own perspective \pc{nowell2017}{3}{A reflexive journal can be used by researchers to record to document the daily logistics of the research}.

\cf{Exemplar study} A five-year mixed-methods case study of nine strategic clinical networks in Alberta. Phase 1 data: 71 documents, 117 interview transcripts and 15 observation field-note sets. Initial codes were generated \emph{deductively} from a pilot, prior research and a conceptual framework, as a variable-oriented strategy for cross-case analysis \pc{nowell2017}{4}{Initial codes were generated deductively based on our pilot study, prior research, and conceptual framework}.

\cf{Phase 1: familiarise} Prolonged engagement, triangulating data modes, documenting reflective thoughts and potential codes, archiving raw data, reading the whole data set at least once before coding \pc{nowell2017}{5}{Braun and Clarke (2006) recommended that researchers read through the entire data set at least once before beginning coding}. Practice: files named by case, source and date, a central archive, and an Excel log of data collection and conversion; heterogeneous documents (Word, PowerPoint, Excel, PDF) needed formatting and OCR \pc{nowell2017}{5}{All files (i.e., raw data) were named to represent the case (i.e., SCN) from which the data came from}.

\cf{Phase 2: initial codes} A coding framework or code manual gives a trail of evidence. Codes need explicit boundaries so they are not interchangeable or redundant, and too many levels of coding hurt clarity \pc{nowell2017}{6}{Codes should have quite explicit boundaries, ensuring they are not interchangeable or redundant}. Apply the chosen technique consistently, document changes in the analytic approach in auditable notations, and re-code data handled the old way \pc{nowell2017}{6}{Any changes to the analytic approach need to be documented in auditable notations}. Credibility is enhanced if more than one researcher analyses the data. Practice: a code manual with detailed definitions and exemplar text (Table 2), two researchers per data set, memos, biweekly peer debriefing with minutes \pc{nowell2017}{7}{The credibility of analysis was further enhanced by having two researchers analyze each data set}.

\cf{Phase 3: searching for themes} Theme definition (DeSantis \& Ugarriza); inductive (data-driven) vs.\ deductive (analyst-driven, less rich) themes, and the need to say which. King's advice: start with a few predefined codes, avoid both too many and too few. Keep a temporary ``miscellaneous'' theme, and do not abandon odd codes \pc{nowell2017}{8}{Braun and Clarke (2006) recommended the creation of a “miscellaneous” theme to temporarily house the codes that do not seem to fit}. Practice: parent nodes from the framework (deductive), child nodes inductive, and diagramming (``around the table'') to explore connections \pc{nowell2017}{9}{We also used diagramming as a way to make sense of the connections between themes}.

\cf{Phase 4: reviewing themes} Check each theme's extracts for coherence and test referential adequacy against the raw data. Add, delete, merge or split codes and themes; some themes lack enough data \pc{nowell2017}{9}{If the researcher identifies a relevant issue in the text not covered by an existing code, a new code may be inserted}.

\cf{Phase 5: defining and naming themes} Detailed written analysis per theme; theme names that convey meaning; the stopping rule that themes are final only when relevant text is no longer left out and the coding has been scrutinised at least twice; team consensus, peer debriefing and written records \pc{nowell2017}{10}{King (2004) suggested that if there remain any sections of text which are clearly relevant to the research question}.

\cf{Phase 6: report} A concise, coherent, non-repetitive account, quotes with identifiers, discrepant data included, process described in sufficient detail, member checking. The authors used the COREQ checklist and only partially member-checked \pc{nowell2017}{11}{We used the Consolidated criteria for reporting qualitative research (COREQ) reporting guidelines}.

\cf{Use for this thesis} A documentation scheme to adopt. Keep an \emph{audit trail} (decision log entries with date, rationale and evidence; versioned codebook with change notes; raw-capture archive with naming convention case--source--date), a reflexive journal, and written minutes of coder discussions. Use Table 1 as a self-audit checklist per phase and state which trustworthiness criteria the thesis addresses and which it does not (e.g.\ no member checking). The paper is about thematic analysis of participant data, so use it for process transparency, not for the walkthrough's observation scheme. It does not give a reliability statistic.
\end{papercard}
